% Book pp. 90-97 (PDF pp. 91-98)
\section{Finding Factors and Zeros of Polynomial Functions}

\begin{activity}{Activity 3.1: Revision of finding quadratic factors}
In pairs or small groups or individually, perform these activities.
\begin{enumerate}
  \item Factorise the following quadratic expressions:
  \begin{enumerate}
    \item[a.] $y^2 + 3y + 2$
    \item[b.] $3y^2 + 11y + 10$
    \item[c.] $6x^2 + 5x + 1$
    \item[d.] $x^2 - 9$
  \end{enumerate}
  \item Write down the factorising process clearly, showing each step taken to
  arrive at the factors.
  \item Exchange your work with other groups or peers to compare answers.
  \item Discuss any differences in methods used and validate the correctness of
  each group's results.
\end{enumerate}
\end{activity}

An expression $(\boldsymbol{x + y})$ is a factor if you can multiply it by another
expression $\boldsymbol{(e + f)}$ to get the original expression
$(\boldsymbol{ex + fx + ey + fy})$. In simple terms, if you can find two or more
expressions that multiply together to make a certain expression, then those are
its factors.

For example, $(x + 2)(x - 3)$ are factors of $x^2 - x - 6$ because expanding
$(x + 2)(x - 3)$:

$x(x - 3) + 2(x - 3)$

$x^2 - 3x + 2x - 6$

$x^2 - x - 6$

Therefore, we see that $(x + 2)(x - 3)$ is indeed equal to $x^2 - x - 6$ confirming
that they are factors of the expression. If you found activity 3.1 challenging, use
the task below to help you.

To \textbf{factorise quadratics} of the form $ax^2 + bx + c$:

\medskip
\noindent\textbf{Example}
\begin{center}
\renewcommand{\arraystretch}{1.3}
\begin{tabular}{|p{0.6\linewidth}|p{0.32\linewidth}|}
\hline
\textbf{Task} & $\boldsymbol{2x^2 + 7x + 6}$ \\
\hline
\textbf{Task 1}\newline
Find two numbers that multiply to give ac (the product of a and c) and add to give b.
& $a = 2$, $b = 7$ and $c = 6$\newline
$ac = 2 \times 6 = 12$\newline
$4 + 3 = 7$ and $4 \times 3 = 12$ \\
\hline
\textbf{Task 2}\newline
Rewrite the middle term using these two numbers as coefficients of x.
& $2x^2 + 4x + 3x + 6$ \\
\hline
\textbf{Task 3}\newline
If the polynomial has four terms, group the terms in pairs.
& $(2x^2 + 4x) + (3x + 6)$ \\
\hline
\textbf{Task 4}\newline
Factorise out the common factor from each group.
& $2x(x + 2) + 3(x + 2)$ \\
\hline
\textbf{Task 5}\newline
Look for a common binomial factor and factorise it out and you have the
factorised solution.
& $(2x + 3)(x + 2)$ \\
\hline
\end{tabular}
\end{center}

\subsection{Factor Theorem and Remainder Theorem}

In addition to quadratic functions, there are several other types of polynomial
functions that can be factorised using the Factor Theorem and the Remainder
Theorem. Remember that if $(x \pm h)$ is a factor of a polynomial, then the remainder
will be zero. Conversely, if the remainder is zero, then $(x \pm h)$ is a factor.

Here are some common types of polynomials:

\subsection{Types of Polynomials}
\begin{enumerate}
  \item \textbf{Cubic Polynomials:} Polynomials of \textit{degree} three, in the form:
  $f(x) = ax^3 + bx^2 + cx + d$
  \item \textbf{Quartic Polynomials:} Polynomials of degree four, in the form:\\
  $f(x) = ax^4 + bx^3 + cx^2 + dx + e$
  \item \textbf{Quintic Polynomials:} Polynomials of degree five, in the form:\\
  $f(x) = ax^5 + bx^4 + cx^3 + dx^2 + ex + f$
\end{enumerate}

\begin{activity}{Activity 3.2: Revision on factorisation}
In groups or individually, use the idea of the factor and remainder theorems to
find the factors or the remainders of the following polynomials.
\begin{enumerate}
  \item Find the remainder when:
  \begin{enumerate}
    \item[a.] $f(x) = 3x^3 - 4x^2 + 2x + 3$ is divided by $(x + 1)$
    \item[b.] $f(x) = 8x^3 - 3x^2 + 6x + 4$ is divided by $(x - 2)$
    \item[c.] $f(x) = 7x^3 + 3x^2 - x + 6$ is divided by $(x + 3)$
    \item[d.] $f(x) = x^3 - x^2 + x + 1$ is divided by $(x - 1)$
  \end{enumerate}
  \item
  \begin{enumerate}
    \item[a.] show that $x + 2$ is a factor of $f(x) = x^3 - 2x^2 - 5x + 6$
    \item[b.] show that $x - 3$ is a factor of $f(x) = x^3 - 3x^2 - 4x + 12$
  \end{enumerate}
  \item Compare your answers with other groups.
\end{enumerate}
\end{activity}

\subsection{Zero-Product Property}

This property of real numbers says that if you multiply two numbers together
and the result is zero, then at least one of those numbers must be zero. In simpler
terms, if you have two numbers, \textbf{a} and b, and when you multiply them
(a$\times$b), and you get zero, it means that either $a$ is zero, or $b$ is zero or
both are zero. This property helps us to find the \textbf{zeros} or \textbf{roots} of
functions of quadratic functions.

The \textbf{zeros} or \textbf{roots} of a function are the values of the input (usually
represented as x) that make the function equal to zero. In other words, if you have
a function f(x), the zeros are the values of x for which f(x) $=$ 0. For example,
given the function $f(x) = 2x^2 - x - 6$, You:
\begin{enumerate}
  \item equate f(x) $=$ 0

  $2x^2 - x - 6 = 0$
  \item factorise the function

  $2x^2 - 4x + 3x - 6 = 0$

  $2x(x - 2) + 3(x - 2) = 0$

  $(2x + 3)(x - 2) = 0$

  Either $(2x + 3) = 0$ or $(x - 2) = 0$
  \item find the solution for each factor

  $(2x + 3) = 0$

  $2x = -3$

  $x = -\frac{3}{2}$

  $(x - 2) = 0$

  $x = 0 - 2$

  $x = 2$
\end{enumerate}
Therefore, the zeros, or roots, of the function $f(x) = 2x^2 - x - 6$ are
$x = -\frac{3}{2}$ and $x = 2$ because both values make the function equal to zero.

We will apply the Zero-Product Property to solve polynomial equations in the
next task.

\subsection{Solving a Polynomial Equation}

There are a number of ways that can be used to find the zeros, or the roots, of
polynomial functions. We can use a calculator, plot a graph, use GeoGebra
(software) and we can also calculate by using the zero property.

\textit{To find the zeros, or roots, by the calculation method:}

Step 1: use the factor theorem approach to factorise the expression completely.

Step 2: equate each factor to 0 and solve.

In pairs or in groups, let's work through the following examples using the steps
above.

\begin{example}{Example 3.1}
Find the zeros of the cubic function $f(x) = x^3 - 2x^2 - 5x + 6$

\solution
$x^3 - 2x^2 - 5x + 6 = 0$

Possible factors of $6 = \pm 1, \pm 2, \pm 3, \pm 6$

Testing with these factors,

When x $=$ 1,

$f(x) = (1)^3 - 2(1)^2 - 5(1) + 6$

$f(x) = 1 - 2 - 5 + 6$

$f(x) = 0$

$\therefore (x - 1)$ is a factor. If x=1 was not a factor, try another option, eg x=-1.

Using the long division method:
\[
\begin{array}{r@{\;}r}
 & x^2 - x - 6 \\ \cline{2-2}
x - 1\,\big) & x^3 - 2x^2 - 5x + 6 \\
 & \underline{-(x^3 - x^2 + 0x + 0)} \\
 & -x^2 - 5x + 6 \\
 & \underline{-(-x^2 + x - 0)} \\
 & -6x + 6 \\
 & \underline{-(-6x + 6)} \\
 & 0
\end{array}
\]
Factorising, $x^2 - x - 6$

$x^2 - 3x + 2x - 6$

$(x^2 - 3x) + (2x - 6)$

$x(x - 3) + 2(x - 3)$

$(x + 2)(x - 3)$

$\therefore \mathrm{f}(\mathrm{x}) = (x - 1)(x + 2)(x - 3)$

$f(x) = 0$

$(x - 1)(x + 2)(x - 3) = 0$

For $(x - 1) = 0$

$x = 1$

For $(x + 2) = 0$

$x = -2$

For $(x - 3) = 0$

$x = 3$

Hence the zeros are when x $= -2, 1, 3$
\end{example}

\begin{example}{Example 3.2}
Find the roots, or zeros, of the cubic function $f(x) = x^3 + x^2 - 10x + 8$

\solution
$x^3 + x^2 - 10x + 8 = 0$

Possible factors of $8 = \pm 1, \pm 2, \pm 4, \pm 8$

Testing with these factors,

When x $=$ 1,

$f(x) = (1)^3 + (1)^2 - 10(1) + 8$

$f(x) = 1 + 1 - 10 + 8$

$f(x) = 0$

$\therefore (x - 1)$ is a factor.

Using the long division method:
\[
\begin{array}{r@{\;}r}
 & x^2 + 2x - 8 \\ \cline{2-2}
x - 1\,\big) & x^3 + x^2 - 10x + 8 \\
 & \underline{-(x^3 - x^2 + 0x + 0)} \\
 & 2x^2 - 10x + 8 \\
 & \underline{-(2x^2 - 2x - 0)} \\
 & -8x + 8 \\
 & \underline{-(-8x + 8)} \\
 & 0
\end{array}
\]
Factorising, $x^2 + 2x - 8$

$x^2 + 4x - 2x - 8$

$(x^2 + 4x) + (2x - 8)$

$x(x + 4) - 2(x + 4)$

$(x + 4)(x - 2)$

$\therefore \mathrm{f}(\mathrm{x}) = (x - 1)(x - 2)(x + 4)$

$f(x) = 0$

$(x - 1)(x - 2)(x + 4) = 0$

For $(x - 1) = 0$

$x = 1$

For $(x - 2) = 0$

$x = 2$

For $(x + 4) = 0$

$x = -4$

Hence the roots, or zeros, are when x $= -4, 2, 1$
\end{example}

\begin{example}{Example 3.3}
Find the zeros of the cubic function. $f(x) = x^3 - 3x^2 - 4x + 12$

\solution
$x^3 - 3x^2 - 4x + 12 = 0$

Possible factors of $12 = \pm 1, \pm 2, \pm 3, \pm 4, \pm 6, \pm 12$

Testing with these factors,

When x $=$ 1,

$f(x) = (1)^3 - 3(1)^2 - 4(1) + 12$

$f(x) = 1 - 3 - 4 + 12$

$f(x) = 6$

$\therefore (x - 1)$ is not a factor and we need to try another potential factor.

When $x = 2$,

$f(x) = (2)^3 - 3(2)^2 - 4(2) + 12$

$f(x) = 8 - 12 - 8 + 12$

$f(x) = 0$

$\therefore (x - 2)$ is a factor.

Using the long division method:
\[
\begin{array}{r@{\;}r}
 & x^2 - x - 6 \\ \cline{2-2}
x - 2\,\big) & x^3 - 3x^2 - 4x + 12 \\
 & \underline{-(x^3 - 2x^2 + 0x + 0)} \\
 & -x^2 - 4x + 12 \\
 & \underline{-(x^2 + 2x - 0)} \\
 & -6x + 12 \\
 & \underline{-(6x + 12)} \\
 & 0
\end{array}
\]
\booknote[S3-01]{the book's long division subtracts $(x^2 + 2x - 0)$ and $(6x + 12)$;
the lines being subtracted are $(-x^2 + 2x - 0)$ and $(-6x + 12)$. The quotient
$x^2 - x - 6$ is correct.}

Factorising, $x^2 - x - 6$

$x^2 - 3x + 2x - 6$

$(x^2 - 3x) + (2x - 6)$

$x(x - 3) + 2(x - 3)$

$(x + 2)(x - 3)$

$\therefore \mathrm{f}(\mathrm{x}) = (x - 2)(x + 2)(x - 3)$

$f(x) = 0$

$(x - 2)(x + 2)(x - 3) = 0$

For $(x - 2) = 0$

$x = 2$

For $(x + 2) = 0$

$x = -2$

For $(x - 3) = 0$

$x = 3$

Hence the zeros are when x $= -2, 2, 3$
\end{example}

Following the remainder theorem is the Rational Zero Theorem which helps to
identify exactly which of the factor will give us zero.
